\chapter{嵌入式系统工程方法}
\label{chap:system-engineering-method}

复杂嵌入式产品的主要风险往往位于模块之间：软件假定电源已稳定，硬件假定驱动不会热插拔，Bootloader 与 OTA 对分区状态理解不同。系统工程的目标是把需求、接口、预算、风险和验证证据组织成可追踪结构。

\section{从利益相关方需求到技术需求}

利益相关方需求描述使用目标和约束，系统需求把它们转换为可验证技术指标。良好需求应具有唯一标识、明确条件、单一义务和验证方法。

“系统应快速启动”不可验证；“在指定电源、温度和存储状态下，从上电到核心功能可用的 99.9 百分位不超过 8 秒”才形成工程约束。

需求应避免提前绑定实现，除非实现本身属于外部约束。安全目标、法规、兼容 ABI 和既有硬件接口通常需要明确指定。

\section{系统边界与上下文}

首先绘制系统上下文：用户、云端、供电、传感器、执行器、生产工装、维修工具和外部网络。每条边界标识数据、能量、信任和生命周期。

边界分析可以发现被遗漏的接口，例如断电期间的机械能量、维修模式的认证、外部 USB 反向供电和工厂网络访问。

\section{架构视图}

单一框图无法表达完整架构，至少需要：
\begin{itemize}
  \item 功能视图：系统提供哪些功能和数据流；
  \item 逻辑视图：模块、任务、服务和接口；
  \item 部署视图：功能运行在哪个核、进程、固件和设备；
  \item 时序视图：启动、控制、故障和更新的事件顺序；
  \item 物理视图：电源域、总线、存储和外部连接；
  \item 安全视图：资产、信任边界和权限；
  \item 验证视图：需求由什么测试和证据关闭。
\end{itemize}

这些视图应共享同一术语和接口 ID，避免不同文档描述出相互矛盾的系统。

\section{接口契约}

接口文档不只列函数或引脚，还应说明：数据类型、单位、范围、时序、阻塞、所有权、错误、版本、复位状态和兼容策略。

\begin{table}[htbp]
  \centering
  \caption{跨层接口契约的必要属性}
  \begin{tabular}{p{30mm}p{110mm}}
    \hline
    属性 & 需要回答的问题 \\
    \hline
    数据 & 编码、单位、范围、字节序、对齐和版本是什么？ \\
    时间 & 周期、截止期、超时、时间戳和抖动上限是什么？ \\
    所有权 & 谁创建、修改、释放和恢复资源？ \\
    故障 & 短传输、重试、重复、复位和部分成功如何处理？ \\
    生命周期 & 上电、升级、休眠、解绑和报废时接口是否存在？ \\
    \hline
  \end{tabular}
\end{table}

接口契约应尽量通过 schema、静态断言、生成代码和兼容测试自动执行。

\section{预算驱动设计}

系统应维护可加和或可分配的预算：启动时间、控制延迟、CPU、RAM、Flash、DDR 带宽、功耗、热、存储寿命和网络流量。

预算需要包含裕量和最坏条件。模块优化若只降低自身平均值，却增加另一个共享资源的峰值，可能使系统整体退化。

每个预算应具有责任人、测量方法和发布门槛，并在硬件或需求变化后重新平衡。

\section{状态机与异常路径}

启动、配网、生产、OTA、低功耗和安全恢复都应使用显式状态机描述。状态机需要定义允许事件、guard、动作、超时和持久化边界。

异常路径不是正常流程末尾的“其他错误”。设计评审应逐一分析依赖未就绪、资源耗尽、部分完成、超时、重复请求、重启和版本不兼容。

\section{架构决策记录}

重要决策应使用简短 ADR 记录背景、选项、决定、后果和重新评估条件。例如选择 A/B 而非单槽更新、选择 RPMsg 而非共享寄存器、选择 IIO 而非私有字符设备。

ADR 记录“为什么”，代码和配置记录“怎么做”。没有理由记录的架构会在人员变化后被反复推翻或错误复制。

\section{风险与技术债务}

风险记录至少包含概率、影响、触发信号、缓解措施和责任人。供应商旧内核、生命周期不稳定器件、未验证机械尺寸和单一密钥都属于需要主动管理的风险。

技术债务应说明当前妥协、影响范围、退出条件和验证缺口。仅写 TODO 而没有优先级和触发条件，通常不会得到治理。

\section{追踪与变更影响}

需求、架构元素、接口、代码、测试和发布证据之间应能追踪。追踪不要求所有关系手工录入大型工具，可以通过稳定 ID、机器可读 manifest 和 CI 报告实现。

变更评审应回答：影响哪些接口、预算、故障假设、安全证据、制造文件和兼容版本。修改 Device Tree 节点名称可能影响启动脚本；替换 PMIC 可能影响 Bootloader、Linux、电源时序和生产测试。

\section{评审与阶段门}

推荐设置需求、架构、原理图、PCB、软件集成、设计验证和量产发布评审。评审门不应只检查文档是否存在，而应检查关键问题是否有证据关闭。

例如设计验证完成需要：需求覆盖、未解决缺陷、性能预算、可靠性、法规、安全和生产测试均达到预定门槛。

\section{文档即代码}

正文、接口 schema、Device Tree binding、寄存器描述、测试计划和发布 manifest 应版本化并在 CI 中构建。图表和生成文件应具有明确来源，避免手工复制后与实现分离。

文档中的代码片段应通过抽取测试或与可构建示例保持关联。无法持续验证的示例容易随 API 演进失效。

\section{完成定义}

嵌入式功能的完成不应止于“代码合并”，而应包括：
\begin{itemize}
  \item 需求和接口已更新；
  \item 正常、错误、恢复和升级路径已经实现；
  \item 资源、时序和功耗预算已经测量；
  \item 单元、集成和目标板测试已经通过；
  \item 诊断、日志和现场恢复能力已经验证；
  \item 构建、符号、SBOM 和发布 manifest 可以复现；
  \item 安全、制造和维护影响已经评审。
\end{itemize}

\section{检查清单}

\begin{itemize}
  \item 需求是否量化、可验证并避免不必要实现绑定？
  \item 是否存在一致的功能、部署、时序、物理和安全视图？
  \item 跨层接口是否定义数据、时间、所有权、故障和生命周期？
  \item CPU、内存、带宽、功耗、热和寿命预算是否持续维护？
  \item 重要决策和技术债务是否记录理由与退出条件？
  \item 变更是否执行跨硬件、软件、测试和生产的影响分析？
  \item 完成定义是否包含验证、诊断、发布和维护证据？
\end{itemize}
